% Diagram and notation definitions shared by the new arXiv manuscript.

\newcommand{\bra}[1]{\langle #1|}
\newcommand{\ket}[1]{|#1\rangle}
\newcommand{\bricknetwork}[1]{%
 \begin{tikzpicture}[x=.43cm,y=.45cm,baseline=-.5ex]
  % the read direction is the flow arrow; no background wash
  \foreach \x in {0,1,2,3}{
   \draw[arxiv boundary] (\x,-.2)--(\x,2.25);}
  \foreach \y in {0,1}{
   \foreach \x in {0,2}{
    \draw[arxiv gate] (\x-.12,\y+.17) rectangle (\x+1.12,\y+.38);}
   \draw[arxiv gate] (.88,\y+.66) rectangle (2.12,\y+.87);
   \draw[arxiv gate] (-.4,\y+.66) rectangle (.12,\y+.87);
   \draw[arxiv gate] (2.88,\y+.66) rectangle (3.4,\y+.87);}
  \draw[arxiv identified] (-.4,-.2) rectangle (3.4,2.25);
  \ifnum#1=1
   \draw[arxiv thermal flow] (3.8,.1)--(3.8,1.9);
  \else
   \draw[arxiv spatial flow] (-.15,2.65)--(3.15,2.65);
  \fi
 \end{tikzpicture}}
\newcommand{\splitgate}{%
 \begin{tikzpicture}[tn,tn centered]
  \draw[phys] (0,-.8)--(0,.8) (1.3,-.8)--(1.3,.8);
  \node[mpo,minimum width=2.2em,minimum height=.5em] at (.65,0) {};
 \end{tikzpicture}}
\newcommand{\bondends}{%
 \begin{tikzpicture}[tn,tn centered]
  \draw[phys] (0,-.8)--(0,.8) (2,-.8)--(2,.8);
  \draw[mpored] (0,0)--(2,0);
  \node[mpo] at (0,0) {$A_\gamma$};
  \node[mpo] at (2,0) {$A_\gamma$};
 \end{tikzpicture}}
\newcommand{\torusgluing}{%
 \begin{tikzpicture}[x=1cm,y=1cm,font=\footnotesize]
  \def\c{.28}\def\g{.12}
  \foreach \a in {0,1}{\foreach \b in {0,1}{
   \tnpatch{{\a*(2*\c+\g)}}{{1.9+\b*(\c+\g)}}{2}{1}{\c}}}
  \foreach \b in {0,1}{
   \tnpatch{5.1}{{1.9+\b*(\c+\g)}}{4}{1}{\c}}
  \foreach \a in {0,1}{
   \tnpatch{{\a*(2*\c+\g)}}{.1}{2}{2}{\c}}
  \tnpatch{5.1}{.1}{4}{2}{\c}
  \node[anchor=north] at (.62,1.85) {$Z_n(\beta)^4$};
  \node[anchor=north] at (5.66,1.85) {$Z_{2n}(\beta)^2$};
  \node[anchor=north] at (.62,.05) {$Z_n(2\beta)^2$};
  \node[anchor=north] at (5.66,.05) {$Z_{2n}(2\beta)$};
  \draw[arxiv spatial flow] (1.45,2.2)--(4.9,2.2)
   node[midway,above,text=black] {$\times\Ps(n,\beta)^2$};
  \draw[arxiv spatial flow] (1.45,.38)--(4.9,.38)
   node[midway,above,text=black] {$\times\Ps(n,2\beta)$};
  \draw[arxiv thermal flow] (.62,1.48)--(.62,.8)
   node[midway,right,text=black] {$\times\Pt(n,\beta)^2$};
  \draw[arxiv thermal flow] (5.66,1.48)--(5.66,.8)
   node[midway,left,text=black] {$\times\Pt(2n,\beta)$};
 \end{tikzpicture}}


% Periodic MPS norm (#1=0) and one-bond energy (#1=1): the ket row and its
% conjugate row each close into a virtual loop drawn clear of the tensors.
\newcommand{\seedmps}[1]{\tikz[tn,tn centered]{
 \foreach \y/\s in {.9/1,-.9/-1}{
  \draw[ring] (-1.1,\y)--(-1.1,{\y+\s*.55})
   --(6.7,{\y+\s*.55})--(6.7,\y)--cycle;}
 \foreach \x in {0,2.2,5.6}{
  \draw[phys] (\x,-.9)--(\x,.9);}
 \tndots{3.9}{.9}\tndots{3.9}{-.9}
 \ifnum#1=1
  \node[sym,fill=mpoblue!14,minimum width=3.4em,minimum height=.65em,inner sep=0pt]
   at (1.1,0) {$h$};
 \fi
 \foreach \x in {0,2.2,5.6}{
  \node[mpo,fill=white] at (\x,.9) {$A$};
  \node[mpo,fill=white] at (\x,-.9) {$A$};}}}


% A reflection-positive bond term with explicit factors: #1 on the left site,
% its mirror image #2 on the right site, across the mirror cut (F = I when
% the two labels agree).
\newcommand{\rpbondpic}[2]{\tikz[tn,tn centered]{
 \draw[phys] (0,-.85)--(0,.85);
 \draw[phys] (3.4,-.85)--(3.4,.85);
 \draw[arxiv bond] (0,0)--(3.4,0);
 \draw[arxiv mirror] (1.7,-1)--(1.7,1);
 \node[xop] at (0,0) {$#1$};
 \node[xop] at (3.4,0) {$#2$};}}
% The bond mirror of the cell theorem: C_alpha against its reflection F C_alpha F.
\newcommand{\scopebond}{\rpbondpic{C_\alpha}{FC_\alpha F}}


\newcommand{\smI}{I}
\newcommand{\smF}{\mathcal F}
\newcommand{\smZ}[1]{Z_{#1}^{(M)}}
\newcommand{\smT}{T_{\beta,M}}
\newcommand{\smD}{\mathcal D_M}
\newcommand{\smnorm}[1]{\left\lVert #1\right\rVert}


\newcommand{\obccap}[2]{%
 \tnchain[capl={\bar v_\beta},capr={#2}]{#1}}
\newcommand{\obcW}{%
 \tnchain[node=mpo,legs,capl={\bar v_\beta},capr={v_\beta}]
 {W_\beta,\cdots,W_\beta}}



% Six-site ring glyph, with a spatial twist on one closing bond when #1=1.
\newcommand{\seedring}[1]{\tikz[baseline=-.6ex,x=1em,y=1em]{
 \draw[black!70,line width=.45pt] (90:1)--(30:1)--(-30:1)
  --(-90:1)--(-150:1)--(150:1)--cycle;
 \ifnum#1=1
  \draw[mpored,line width=1.6pt] (150:1)--(90:1);
  \node[overlay,font=\scriptsize,text=mpored] at (-.8,1.15) {$\pi$};
 \fi
 \foreach \a in {90,30,-30,-90,-150,150}{
  \filldraw[draw=black!70,fill=white,line width=.45pt] (\a:1) circle (.12);}}}

% Extremal twelfth-moment weights and their sixth powers, not an observed spectrum.
\newcommand{\seedweightbars}{\begin{tikzpicture}[x=1cm,y=.55cm,baseline=-.5ex]
 % half-height bars with the labels beside them, so the glyph costs three text lines
 \fill[arxiv leading weight] (0,0) rectangle (.4,.9881);
 \fill[arxiv remainder weight] (.8,0) rectangle (1.2,.0777);
 \draw[arxiv guide] (-.1,.9881)--(1.3,.9881);
 \node[arxiv label,anchor=west,inner sep=2pt] at (1.3,.9881) {$\ell^{12}$};
 \node[arxiv label,anchor=west,inner sep=2pt] at (1.3,.12) {$m_J-\ell^{12}$};
 \draw[arxiv gray flow] (3.1,.5)--(4.1,.5)
  node[midway,above,arxiv label,inner sep=1.5pt] {$u\mapsto u^6$};
 \fill[arxiv leading weight] (4.6,0) rectangle (5,.9305);
 \fill[arxiv remainder weight] (5.4,0) rectangle (5.8,.000001);
 \node[arxiv label,anchor=west,inner sep=2pt] at (5.9,.9305) {$\ell^{72}$};
 \draw[arxiv baseline] (-.1,0)--(1.3,0) (4.5,0)--(5.9,0);
\end{tikzpicture}}

% The four established coupling/squaring steps, in the talk's two-lane layout.
% Rows are packed to 1.05 cm, and the thermal lane sits half a row above
% the spatial one, so the staircase fits a column without a tall gap.
\newcommand{\pbcsteps}{\begin{prooflanes}\def\lanedy{1.05}
 \lanebound{s}{ps0}{0}{0}{$\Ps(n,\beta)\ge1-p$}
 \lanebound{t}{pt0}{1}{0}{$\Pt(2n,\beta)\ge1-q$}
 \lanebound{so}{ps1}{0}{1}{$\Ps(n,2\beta)\ge(1-p)^2(1-q)$}
 \lanebound{s}{ps2}{0}{2}{$\Ps(2n,2\beta)\ge1-p'$}
 \lanebound{to}{pt1}{1}{2.5}{$\Pt(4n,\beta)\ge(1-q)^2(1-p')$}
 \lanebound{t}{pt2}{1}{3.5}{$\Pt(4n,2\beta)\ge1-q'$}
 \lanemove{ps0}{ps1}{1}{identity}\laneuse{pt0}{1}
 \lanemove[mpored]{ps1}{ps2}{2}{square in space}
 \lanemove{pt0}{pt1}{3}{identity}\laneuse{ps2}{3}
 \lanemove[mpoblue]{pt1}{pt2}{4}{square in time}
\end{prooflanes}}

% Adapted from the proof-strategy slide; errors are schematic, not the seed orbit.
\newcommand{\strategystairs}{\begin{tikzpicture}[x=.8cm,y=.8cm,
 every node/.style={arxiv label,inner sep=1pt}]
 \draw[arxiv dotted guide] (0,0)--(3.3,2.1);
 \foreach \x/\y in {0/0,1.65/1.05}{
  \draw[arxiv thermal flow,line width=1pt] (\x,\y+.04)--(\x,\y+1.01);
  \draw[arxiv spatial flow,line width=1pt,densely dashed]
   (\x+.04,\y+1.05)--(\x+1.61,\y+1.05);}
 \foreach \x/\y in {0/0,0/1.05,1.65/1.05,1.65/2.1,3.3/2.1}{
  \node[arxiv point] at (\x,\y) {};}
 \foreach \x/\lab in {0/n_0,1.65/2n_0,3.3/4n_0}{
  \node[anchor=north] at (\x,-.16) {$\lab$};}
 \foreach \y/\lab in {0/\beta_0,1.05/2\beta_0,2.1/4\beta_0}{
  \node[anchor=east] at (-.15,\y) {$\lab$};}
 \node[anchor=west,fill=white] at (.12,.22) {$\varepsilon$};
 \node[anchor=west,fill=white] at (1.8,1.28) {$\approx\varepsilon^2$};
 \node[anchor=south] at (3.3,2.22) {$\approx\varepsilon^4$};
 % key: each move with the arrow that draws it, close to the staircase
 \draw[arxiv thermal flow,line width=1pt] (4.15,1.9)--(4.15,2.35);
 \node[text=mpoblue,anchor=west] at (4.3,2.12) {double $\beta$};
 \node[anchor=west] at (4.3,1.68) {$\rho_n(\beta)^2\propto\rho_n(2\beta)$};
 \draw[arxiv spatial flow,line width=1pt,densely dashed] (3.95,.95)--(4.4,.95);
 \node[text=mpored,anchor=west] at (4.5,.95) {double $n$};
 \node[anchor=west] at (4.3,.5) {$\sigma_n^2\propto\sigma_{2n}$};
\end{tikzpicture}}

% ------------------------------------------------- Supplemental Material glyphs
% One site of the brickwork network through three slices, read along space:
% the odd-bond labels q_t enter from the left and the even-bond labels p_t
% leave to the right; the loop is the physical trace that gives F_{pq}.
\newcommand{\sitetrace}{\tikz[tn,tn centered,font=\scriptsize]{
 \draw[phys,ring] (0,-1.75)--(0,1.75)--(2.9,1.75)--(2.9,-1.75)--cycle;
 \foreach \y/\p/\q in {-1.05/p_0/q_0,0/p_1/q_1,1.05/p_2/q_2}{
  \draw[arxiv bond] (-1.05,{\y-.22})--(0,{\y-.22});
  \draw[arxiv bond] (0,{\y+.22})--(1.05,{\y+.22});
  \node[xvertex] at (0,{\y-.22}) {};
  \node[xvertex] at (0,{\y+.22}) {};
  \node[anchor=east,inner sep=1pt] at (-1.05,{\y-.22}) {$\q$};
  \node[anchor=west,inner sep=1pt] at (1.05,{\y+.22}) {$\p$};}}}

% One site of the brickwork network traced through M=2 slices, read along
% space: on each slice the odd-bond factor A_{a_t} enters from the left and the
% even-bond factor A_{b_t} leaves to the right; the closed physical loop is the
% trace, so the picture is the matrix element F_{ba} of the main text.
\newcommand{\sitetraceab}{\tikz[tn,tn centered,font=\scriptsize]{
 \draw[phys,ring] (0,-1.2)--(0,1.2)--(2.7,1.2)--(2.7,-1.2)--cycle;
 \foreach \y/\b/\a in {-.55/b_1/a_1,.55/b_2/a_2}{
  \draw[arxiv bond] (-1.05,{\y-.22})--(0,{\y-.22});
  \draw[arxiv bond] (0,{\y+.22})--(1.05,{\y+.22});
  \node[xvertex] at (0,{\y-.22}) {};
  \node[xvertex] at (0,{\y+.22}) {};
  \node[anchor=east,inner sep=1pt] at (-1.05,{\y-.22}) {$\a$};
  \node[anchor=west,inner sep=1pt] at (1.05,{\y+.22}) {$\b$};}}}

% \capbars[yunit]{h1,h2,...}{cap}: twelfth-power weights filled up to a cap
% (dashed) with the remaining mass in one smaller entry; the extremal
% configuration of the capped-mass bound.
\newcommand{\capbars}[3][1.1em]{\tikz[baseline=0pt,x=.5em,y=#1]{
 \foreach \h [count=\i from 0] in {#2}{
  \fill[arxiv leading weight] ({1.15*\i},0) rectangle ++(.8,\h);
  \xdef\capn{\i}}
 \draw[arxiv guide] (-.2,#3)--({1.15*\capn+1},#3);
 \draw[arxiv baseline] (-.25,0)--({1.15*\capn+1.05},0);}}

% Twelfth-power weights of K at the base seed: one dominant J eigenvalue,
% the J remainder, and the three identical N lists.
\newcommand{\sectorspectrum}{\tikz[baseline=0pt,x=.34em,y=1.1em]{
 \fill[arxiv leading weight] (0,0) rectangle ++(.78,1);
 \foreach \h [count=\i] in {.35,.2}{\fill[arxiv remainder weight] (\i,0) rectangle ++(.78,\h);}
 \foreach \h [count=\i from 3] in {.58,.58,.58,.3,.3,.3}{\fill[arxiv bulk weight] (\i,0) rectangle ++(.78,\h);}
 \draw[arxiv baseline] (-.2,0)--(9,0);
}}

% The filter F_N against its certified total B_N: the level is crossed at
% +-V and no eigenvalue of the N lists can sit on the shaded exterior rays.
\newcommand{\filterglyph}{\tikz[baseline=0pt,x=1.7em,y=1.1em,line width=.45pt]{
 \clip (-1.05,-.08) rectangle (1.05,1.05);
 \fill[arxiv exclude band] (-1.05,0) rectangle (-.872,1.05);
 \fill[arxiv exclude band] (.872,0) rectangle (1.05,1.05);
 \draw[arxiv baseline] (-1.05,0)--(1.05,0);
 \draw[arxiv guide] (-1.05,.535)--(1.05,.535);
 \draw[arxiv curve,domain=-1:1,samples=161,smooth]
  plot (\x,{(\x^4)*((12-394*\x*\x+1000*\x^4)/300)^2});}}

% The six exact rational steps on a logarithmic scale: data of Table sm:tab:orbit.
\newcommand{\defectladder}{\begin{tikzpicture}[x=1.05cm,y=1.6cm,font=\footnotesize]
 \def\ps{-1.320,-1.218,-1.163,-1.146,-1.198,-1.425,-2.072}
 \def\qs{-.741,-.761,-.787,-.839,-.976,-1.350,-2.258}
 \draw[arxiv guide] (0,-2)--(6,-2)
  node[right,arxiv note] {$10^{-2}$};
 \draw[arxiv dotted guide] (0,-.895)--(6,-.895)
  node[right,arxiv note] {$\varepsilon_*$};
 \foreach \v [count=\i from 0] in \ps{\ifnum\i>0 \draw[arxiv series,mpored] ({\i-1},\pl)--(\i,\v);\fi
  \fill[mpored] (\i,\v) circle[radius=1.3pt]; \xdef\pl{\v}}
 \foreach \v [count=\i from 0] in \qs{\ifnum\i>0 \draw[arxiv series,mpoblue] ({\i-1},\ql)--(\i,\v);\fi
  \fill[mpoblue] (\i,\v) circle[radius=1.3pt]; \xdef\ql{\v}}
 \node[mpored,anchor=north] at (2.5,-1.19) {$p_j$ (spatial)};
 \node[mpoblue,anchor=south] at (1.5,-.76) {$q_j$ (thermal)};
 \foreach \n [count=\i from 0] in {36,72,144,288,576,1152,2304}{
  \node[arxiv note] at (\i,-2.42) {$\n$};}
 \node[arxiv note,anchor=east] at (-.25,-2.42) {$n_j$};
\end{tikzpicture}}

% Interpolation between dyadic scales, adapted from the talk: the unknown
% length L is bounded from above at beta_j by the torus on its left and from
% below at 2 beta_j by the torus on its right.
\newcommand{\dyadicinterval}{\begin{tikzpicture}[x=1cm,y=.75cm,
 every node/.style={arxiv label,inner sep=1pt}]
 \draw[black!25,line width=.45pt] (-.3,0)--(4.4,0) (-.3,1.8)--(4.4,1.8);
 \foreach \x/\l in {0/n_j,3.6/2n_j}{
  \draw[black!25,line width=.45pt] (\x,-.2)--(\x,2.2);
  \node[anchor=north] at (\x,-.25) {$\l$};}
 \node[anchor=east] at (-.35,0) {$\beta_j$};
 \node[anchor=east] at (-.35,1.8) {$2\beta_j$};
 \draw[arxiv dotted guide] (0,0)--(3.6,1.8);
 \draw[tncert,line width=1pt] (1.9,-.2)--(1.9,2.2);
 \node[anchor=north,arxiv cert note] at (1.9,-.25) {$L$};
 \draw[arxiv spatial flow,line width=.7pt] (.12,0)--(1.78,0);
 \draw[arxiv spatial flow,line width=.7pt] (3.48,1.8)--(2.02,1.8);
 \node[arxiv point] at (0,0) {};
 \node[arxiv point] at (3.6,1.8) {};
 \node[arxiv point,fill=tncert] at (1.9,0) {};
 \node[arxiv point,fill=tncert] at (1.9,1.8) {};
 \node[anchor=south,text=mpored,inner sep=1pt,fill=white] at (.95,.07) {power $L/n_j$};
 \node[anchor=south west,text=mpored,inner sep=1pt,fill=white] at (2.05,1.86) {power $2n_j/L$};
 \node[anchor=west,arxiv note] at (2.45,.62) {doubling};
\end{tikzpicture}}

% Word-sized dyadic interval: the two known tori (gray points, at (n_j, beta_j)
% below and (2n_j, 2beta_j) above) bound the thermal purity at the target length
% L (teal) from either side by the power-sum inequality along the length axis
% (red). Unlabeled so that it stays one line tall; the prose names the points.
\newcommand{\interpolationglyph}{\tikz[baseline=-.1em,x=1em,y=.8em]{
 \draw[arxiv baseline] (-.15,0)--(2.75,0) (-.15,.75)--(2.75,.75);
 \draw[arxiv series,tncert] (1.3,-.1)--(1.3,.87);
 \draw[arxiv small arrow,mpored] (.2,0)--(1.18,0);
 \draw[arxiv small arrow,mpored] (2.4,.75)--(1.42,.75);
 \node[arxiv point,inner sep=.9pt] at (0,0) {};
 \node[arxiv point,inner sep=.9pt] at (2.6,.75) {};
 \node[arxiv point,inner sep=.9pt,fill=tncert] at (1.3,0) {};
 \node[arxiv point,inner sep=.9pt,fill=tncert] at (1.3,.75) {};
}}

